\documentclass[11pt]{article}

\setlength{\parindent}{0pt}
\usepackage{xltxtra}
\usepackage{hyperref}
\hypersetup{hidelinks}
\usepackage{url}
\urlstyle{tt}
\usepackage{xcolor}
\definecolor{CVBlue}{RGB}{23,110,191}
\usepackage{calc}
\usepackage{graphicx}
\usepackage{tikz}
\usetikzlibrary{calc}
\usepackage{fontspec}
\usepackage{xeCJK}
\usepackage{enumitem}
\CJKsetecglue{}

%% 主文档字体设置：保持原模板
\setmainfont[
    Path = fonts/Main/,
    Extension = .otf,
    BoldFont = texgyretermes-bold.otf,
]{texgyretermes-regular.otf}

%% 中文字体设置：保持原模板
\setCJKmainfont[
    Path = fonts/hansans/,
    Extension = .ttf,
    BoldFont = NotoSansSC-Bold.ttf,
]{NotoSansSC-Regular.otf}

\usepackage{relsize}
\usepackage{xspace}
\usepackage{fontawesome}

\usepackage[
    a4paper,
    left=1.2cm,
    right=1.2cm,
    top=1.5cm,
    bottom=1cm,
    nohead
]{geometry}

\renewcommand{\baselinestretch}{1.5}

\usepackage{titlesec}
\usepackage{enumitem}
\setlist{noitemsep}
\setlist[itemize]{topsep=0.25em, leftmargin=*}
\setlist[enumerate]{topsep=0.25em, leftmargin=*}

\newlength{\interProjectSpacing}
\setlength{\interProjectSpacing}{0.9em}
\newcommand{\projectsep}{\vspace{\interProjectSpacing}}

\newlength{\intraProjectTitleSep}
\setlength{\intraProjectTitleSep}{0.4em}
\newcommand{\titlebreak}{\\[\intraProjectTitleSep]}

\newlength{\intraProjectListTopSep}
\setlength{\intraProjectListTopSep}{0.2em}

\titleformat{\section}
{\large\bfseries\raggedright}
{}{0em}
{}
[{\color{CVBlue}\titlerule}]
\titlespacing*{\section}{0cm}{*1.6}{*1.2}

\begin{document}
\pagenumbering{gobble}

%%%% 照片：保留，文件名为 avatar.jpg
\begin{tikzpicture}[remember picture, overlay]
    \node[anchor = north east] at ($(current page.north east)+(-2cm,-1.2cm)$)
    {\includegraphics[height=3cm]{avatar.jpg}};
\end{tikzpicture}

%%%% 浙大 Logo：保留，文件名为 zju.png
\begin{tikzpicture}[remember picture, overlay]
    \node[anchor = north west] at ($(current page.north west)+(0.5cm,+1.0cm)$)
    {\includegraphics[height=6cm]{zju.png}};
\end{tikzpicture}

\centerline{\LARGE\bfseries{朱泽豪}}

\centerline{\normalsize{\faPhone\ 13676835081 \quad
\faEnvelopeO\ \href{mailto:3240103239@zju.edu.cn}{3240103239@zju.edu.cn}}}

\centerline{\normalsize{\faRssSquare\ \href{https://www.cnblogs.com/Plozia}{www.cnblogs.com/Plozia}}}

\section{\makebox[\widthof{\faGraduationCap}][c]{\color{CVBlue}\faGraduationCap}\ 教育背景}
\textbf{浙江大学 \quad 数学科学学院 \quad 信息与计算科学} \hfill 2024.9 -- 至今\\[0.5em]
大三在读
\begin{itemize}[nosep]
    \item 相关课程：《数值代数》《数值分析》《数学建模（H）》《高级数据结构与算法分析》《数学实践》
\end{itemize}

\section{\makebox[\widthof{\faUsers}][c]{\color{CVBlue}\faUsers}\ 竞赛经历}

% --- 美赛 ---
\textbf{2026 美国大学生数学建模竞赛（MCM/ICM）\quad H 奖} \hfill 2026.2 \titlebreak
项目描述：参与预测/优化模型的编程实现，使用 Python/C++ 完成模型求解、结果分析与可视化；结合 AI 辅助完成代码实现与模型迭代。
\begin{itemize}[nosep, topsep=\intraProjectListTopSep]
    \item \textbf{编程实现}：使用 \textbf{Python/C++} 完成模型预测、微分方程求解与结果分析，结合 AI 辅助完成代码实现与模型迭代。
    \item \textbf{优化与模拟}：围绕物资运输方案构建多目标优化模型，求解帕累托解集；引入 \textbf{蒙特卡洛模拟}分析故障对成本与时间的影响，并完成单变量敏感性分析与结果可视化。
    \item \textbf{竞赛成果}：美国大学生数学建模竞赛 \textbf{H 奖}。
\end{itemize}

\projectsep

% --- 国赛 ---
\textbf{2026 全国大学生数学建模竞赛\quad 国家一等奖} \hfill 2026.9 \titlebreak
项目描述：针对中药材热风烘干过程，建立一维径向传热传质耦合模型，负责给定微分方程的数值求解与程序实现。
\begin{itemize}[nosep, topsep=\intraProjectListTopSep]
    \item \textbf{数值离散}：采用守恒\textbf{有限体积法}进行空间离散，使用 Rannacher 启动的 \textbf{Crank-Nicolson/后向 Euler} 时间推进，处理表面 Robin 边界与中心对称条件。
    \item \textbf{非线性与线性求解}：通过 \textbf{Picard/不动点迭代}处理变物性与非线性扩散系数，使用三对角/带状 LU 求解线性系统；参与移动边界模型的坐标变换与数值实现。
    \item \textbf{竞赛成果}：完成温度场与水分场计算及干燥时长预测，获全国大学生数学建模竞赛 \textbf{国家一等奖}。
\end{itemize}

\section{\makebox[\widthof{\faUsers}][c]{\color{CVBlue}\faUsers}\ 校园工作}
\textbf{浙江大学学生无极棋社\quad 人资部干事} \hfill 2024.9 -- 2025.9 \titlebreak
任职期间考核良好，参与社团人力资源相关工作。

\section{\makebox[\widthof{\faCogs}][c]{\color{CVBlue}\faCogs}\ 技术栈}
\begin{itemize}[nosep]
    \item \textbf{编程语言：} Python, C++
    \item \textbf{工具与辅助：} AI 辅助编程、LaTeX
    \item \textbf{英语水平：} CET-6 511
\end{itemize}

\section{\makebox[\widthof{\faGraduationCap}][c]{\color{CVBlue}\faList}\ 获奖情况}
\begin{itemize}
    \item 全国大学生数学建模竞赛 \quad 国家一等奖 \hfill 2026.9
    \item 美国大学生数学建模竞赛 \quad H 奖 \hfill 2026.2
\end{itemize}

\section{\makebox[\widthof{\faInfo}][c]{\color{CVBlue}\faInfo}\ 其他}
\begin{itemize}[parsep=0.5ex]
    \item \textbf{个人博客：} \href{https://www.cnblogs.com/Plozia}{www.cnblogs.com/Plozia}（高中时期维护，大学未持续更新）
\end{itemize}

\end{document}
